\section{Limitations and threats to validity}
\label{sec:limitations}

A machine-checked proof removes one risk, an error in the argument. It
leaves the others in place. We list them from the most to the least
important.

\paragraph{The Dex side is a model, not the code.}
The semantics is written by hand from the server's source and validated by
review. It is not extracted from the implementation, and no tool checks
that the server refines it. A mistake in reading the code is a mistake in
the theorem's meaning. Three things reduce the risk and none removes it.
The rules cite the implementing code, except one RPC rule that follows a
fix not yet merged at the pinned commit. Three adversarial reviews checked
the rules against that code. And the model's executable form runs the
examples and the lockstep checks, although only against itself, not against
a running server. The Lean model is, further, a \emph{restriction} of the
specification to what the translation and the admissible environment use;
each omission is listed in the development with the reason why it cannot
matter. The specification also pins a base commit of the server. The server
has changed since, including the two fixes of
\S\ref{sec:server-bugs}; the second changes what happens to the messages of
a canceled Step, a behavior that a translated Flow reaches only at its
close, where it is unobservable. That last sentence is an argument, not a
theorem. The failure rules cite a later commit of the server and went
through their own adversarial reviews, with the same caveat.

\paragraph{Failures: scope and hypotheses.}
The failure theorem of \S\ref{sec:failures} replaces the failure-free
assumption of the first theorem with a list of narrower ones. Each is a
place where a real deployment can leave the theorem.
\begin{itemize}
\item \emph{The Lean fragment.} One Flow, a start and publications, and
  StartFlow retries. RPCs, SubFlows, StopFlow, and Attribute locks under
  failures are not covered.
\item \emph{SYNC only.} An ASYNC Step method can lose a start in a crash
  and run the same attempt again with a new snapshot, so its effects can
  differ. That is the real ASYNC behavior, and the theorem excludes it.
\item \emph{H1 to H4.} No method exhausts its retries, no Step uses a
  failure policy, user code ignores the attempt number, the first attempt
  time, the last error, and the last heartbeat value, and no failure is
  non-retryable. A long outage can break H1: an interrupt after an outage
  longer than the total retry duration exhausts the method. Code that
  resumes from a heartbeat checkpoint breaks H3. The Lean model leaves
  failure policies and ASYNC starts out, so H2 and SYNC hold there by
  construction; the theorem covers those features by not modeling them.
\item \emph{Finitely many failures for liveness.} Liveness also assumes
  that the server comes back after every crash and that every interrupt
  commits. Progress is defined through the abstraction map, so the theorem
  says nothing about an internal livelock of the engine.
\item \emph{The clock.} On one node the wall clock must not step back
  between the timer service's due check and the stamping of its command.
  The theorem speaks of stamps: a host clock that runs ahead can fire a
  timer early in real time.
\item \emph{Publications retried by the client.} A publication has no
  request ID. A client that sends it again after a lost response has sent a
  second request, which the server accepts as a second message. That is
  outside failure transparency. Exactly-once publication needs another
  route, such as an RPC with a caller-chosen request ID.
\item \emph{StartFlow retries after the purge.} A retry deduplicates only
  while the run is still in the Flow directory; after the purge it starts a
  new run. The theorem assumes every retry comes before the purge, and the
  Lean model, which keeps a closed run forever, makes this hold by
  construction.
\item \emph{Liveness backward only without retries.} Every fair DexF trace
  erases to a fair Dex trace, but the converse holds only for input words
  without StartFlow retries, because stamped retries move the clock
  (\S\ref{sec:failure-theorems}).
\item \emph{Timer order is idealized.} Both theorems fire timers as a
  prefix of the due timers. The server can fire a timer while an earlier-due
  timer, armed by a command applied just before the scan's command, is
  still armed. The owner kept the prefix rule, because no distributed
  system guarantees the order of timers due at almost the same time.
  Neither theorem covers runs in that window.
\item \emph{One owner at a time.} A standby takeover is in scope.
  Multi-node partition ownership is not: the server runs one node per
  cluster, and the model has one log writer at a time.
\end{itemize}
The theorem also says less than a reader might hope about effects.
Repeated effects may overlap in time and come late: a zombie can emit an
effect after a newer attempt of the same method has finished. So an
idempotent effect is not enough for exactly-once in effect. Idempotent
does not mean commutative: a zombie's late \emph{set $x$ to 1} can
overwrite a later \emph{set $x$ to 2}, although each call is idempotent.
Exactly-once in effect needs idempotency keys, commutative effects, or
effects fenced by attempt. The server's deterministic simulation harness
with injected crashes checks that a recovered store equals one that never
crashed; it is complementary evidence and is not connected to the model.

\paragraph{Fairness is assumed, not provided.}
Liveness holds for fair runs, and the fairness conditions of
Table~\ref{tab:fairness} are assumptions about components outside the
rules. F2 assumes that the timer service eventually fires a due timer that
nothing else disarms. F3 assumes that the dispatcher eventually starts a
pending method that is not canceled first. F5 assumes
that the engine eventually commits or drops a returned response. F4 is the
one most worth stating plainly: it assumes that every Worker task
eventually settles, by returning or by being abandoned after its invocation
was canceled. The server does not guarantee this on its own. Cancellation
of a running handler is cooperative, and in the failure-free model there is
no method timeout to end a handler that never returns. Under failures a
timeout fails the attempt, but the handler may keep running as a zombie,
so the failure theorem's liveness keeps the same assumption for zombies. A deployment in
which a handler can hang is outside the liveness theorem; the safety
theorem still applies to it. The fair runs that the development exhibits
(\S\ref{sec:nonvacuity}) cover termination, deadlock, and divergence by an
infinite input word. A fair run of a program that diverges on its own, on a
finite word, is not exhibited, and no lemma says that every reachable
configuration has a fair continuation.

\paragraph{A restricted environment.}
The admissible environment starts one Flow and publishes plain data to
external Channels. It makes no RPC, never stops a Flow, never sets an
Attribute from outside, and never deletes a message or skips a timer. The
specification gives rules for all of these, but they are outside the
theorem. In particular an RPC can observe intermediate state and can fail
with a lock conflict, and stopping a Flow cancels Steps at arbitrary
points; both would need \Lang{} counterparts.

\paragraph{Well-formed programs only.}
The hypothesis \lean{Program.IsWellFormed} excludes programs that can break
a mutex requirement. For such a program the two sides differ: \Lang{}
refuses the step, and the Flow takes it. The scope is semantic and
undecidable in general. The static lock discipline of
\S\ref{sec:nonvacuity} is sufficient, not necessary. It tracks statically
known mutexes only, and it empties the lock set after a call and after a
\kw{try}. So it rejects \kw{unlock}, \kw{wait}, \kw{signal}, and
\kw{broadcast} on a created mutex, and a program that holds a lock across a
call and unlocks after it, although such programs can be well formed. For
them well-formedness must be proved by other means.

\paragraph{The timing model.}
\Lang{} counts stamped actions exactly as the translation counts commits
(\S\ref{sec:timing}). The theorem is exact for this model and silent about
any other. A program whose observable behavior depends on the order of two
deadlines that are close together can behave differently under another cost
model. We consider such programs incorrect, but the theorem does not
exclude them; it just assigns them the Dex timing. The equivalence is also
untimed: it does not say that an event happens within some time, only that
the same sequences of events are possible. Durations are whole seconds.

\paragraph{Unbounded counters.}
Step numbers and timer numbers are unbounded in the model and 32-bit in the
server. The order of timers with equal fire times was derived from the
protobuf encoding of timer IDs by hand; that derivation is an assumption of
the model, and it holds for fewer than $2^{32}$ Step executions of one type
per run.

\paragraph{\Lang{} is a chosen language.}
\Lang{}'s semantics was designed so that the equivalence is exact: FIFO
mutexes and condition variables, cooperative cancellation at four
cancellation points, asynchronous \kw{emit}, a \kw{select} with one receive
branch and one timeout branch, late timers, and a single-use join. Each
choice is defensible and close to common practice, but \Lang{} is not an
existing language, and a different reasonable choice (say, a cancellable
\kw{lock}) would need a different translation and a new proof. The theorem
shows that Dex can express \emph{a} conventional concurrent language. It
does not show that every feature of every such language has a faithful
encoding.

\paragraph{Observations.}
Only inputs, external effects, and the program's result are observable.
Reads of a Flow's state by a client are not, and the events of the Flow's
internal Channels are projected away. An observer with access to those
could tell a translated Flow from its source program.

\paragraph{Trust in the statements.}
The kernel checks proofs, not intentions. A reader who wants to rely on the
first theorem must read the model, the translation, and the statements,
about 3,450 lines, against the specification; the failure theorem adds
about 1,900 lines of definitions and statements. The development maps every section
and every numbered assumption of the specification to its Lean definition
to make that reading possible.
